\documentclass{article}

\usepackage[final]{neurips}

\usepackage[utf8]{inputenc}
\usepackage[T1]{fontenc}
\usepackage{hyperref}
\usepackage{url}
\usepackage{booktabs}
\usepackage{amsfonts}
\usepackage{amsmath}
\usepackage{nicefrac}
\usepackage{microtype}
\usepackage{xcolor}
\usepackage{graphicx}
\usepackage{array}
\usepackage{multirow}

% TODO marker
\newcommand{\todo}[1]{\textcolor{red}{[TODO: #1]}}

\title{A Lightweight Fine-Tuned LLM Shopping Assistant with Structured and Semantic Product Retrieval}

\author{%
  \todo{Member 1 Name}\\
  Department of Electrical and Computer Engineering\\
  University of Toronto\\
  \texttt{\todo{email}} \\
  \And
  \todo{Member 2 Name}\\
  Department of Electrical and Computer Engineering\\
  University of Toronto\\
  \texttt{\todo{email}} \\
  \And
  \todo{Member 3 Name}\\
  Department of Electrical and Computer Engineering\\
  University of Toronto\\
  \texttt{\todo{email}} \\
}

\begin{document}

\maketitle

% =========================================================
\begin{abstract}

This project develops a lightweight large language model (LLM)-based shopping
assistant for retrieving products from a large Amazon product database.
Instead of relying entirely on a general-purpose LLM to generate product
recommendations, we decompose the task into several controlled components:
structured query extraction, fine-tuned product category classification,
database filtering, semantic retrieval, and recommendation generation.

A key component of the system is a lightweight approximately one-billion
parameter language model that maps a product request to one of 16 predefined
main product categories. We compare three model families: Qwen3.5-0.8B,
Llama 3.2-1B, and Gemma 3 1B IT. The models are fine-tuned for closed-set
product category classification using parameter-efficient fine-tuning.
Gemma 3 1B IT achieved the best observed classification performance with
93.5\% accuracy and a Macro F1 score of 87.4\%, and was therefore selected
for integration into the final retrieval pipeline.

The final system combines the fine-tuned classifier with structured SQL
filters and semantic product embeddings. This architecture improves control
over product type, price, rating, and other user constraints while keeping
the computational requirements substantially lower than those required by
larger language models.

\end{abstract}


% =========================================================
\section*{Statement of Teamwork}

The project was completed collaboratively by three team members.

\begin{itemize}
    \item \textbf{\todo{Member 1}:}
    \todo{Describe contribution, e.g., Qwen fine-tuning, evaluation,
    dataset preprocessing, presentation.}

    \item \textbf{\todo{Member 2}:}
    \todo{Describe contribution, e.g., Llama fine-tuning, vector database,
    retrieval system, integration.}

    \item \textbf{\todo{Member 3}:}
    Implemented and evaluated the Gemma 3 1B IT fine-tuning pipeline,
    performed LoRA-based supervised fine-tuning for product main-category
    classification, evaluated zero-shot/prompt-based and fine-tuned
    performance, and integrated category prediction with the structured
    product search pipeline.
\end{itemize}


% =========================================================
\section{Introduction}

Online shopping platforms contain millions of products, making it difficult
for users to manually search through all available items. Traditional
keyword search can filter products using exact terms, but it may not fully
understand natural-language requests such as:

\begin{quote}
``I want a reliable wireless gaming mouse under \$50 with good reviews.''
\end{quote}

Such a request includes multiple types of information: product type,
price constraints, quality requirements, and semantic preferences.
A useful shopping assistant must understand these different constraints and
translate them into an effective product retrieval query.

One possible solution is to directly use a general-purpose LLM such as
ChatGPT to answer the request. However, this approach introduces several
limitations. In our exploratory testing, when asking a general-purpose model
for products such as ``a mouse under \$50,'' some returned recommendations
could include related accessories such as mouse pads rather than strictly
matching the requested product type. This behavior is understandable for a
generative language model because it is optimized to produce semantically
relevant text rather than perform exact closed-set product classification.

For our application, category precision is important because incorrectly
identifying a product category can cause the retrieval stage to search an
incorrect region of the database. Therefore, instead of asking a large
general-purpose model to perform every step, we use a modular architecture.

The system consists of:

\begin{enumerate}
    \item Natural-language query understanding;
    \item Structured attribute extraction;
    \item Fine-tuned main-category classification;
    \item SQL filtering;
    \item Semantic vector retrieval;
    \item Final recommendation generation.
\end{enumerate}

The main deep learning problem investigated in this project is whether a
small language model, approximately one billion parameters in size, can be
fine-tuned to accurately map product titles or user product requests into a
fixed taxonomy of 16 product categories.

We compare three lightweight model families:

\begin{itemize}
    \item Qwen3.5-0.8B;
    \item Llama 3.2-1B;
    \item Gemma 3 1B IT.
\end{itemize}

The goal is not to benchmark every available small language model. Instead,
these three models were chosen because they represent different widely used
model families while remaining in a similar parameter range. This allows a
more meaningful comparison under similar computational constraints.


% =========================================================
\section{Preliminaries and Problem Formulation}

\subsection{Dataset}

The project uses an Amazon product dataset containing approximately
1.43 million product records.

Each product contains attributes such as:

\begin{itemize}
    \item ASIN;
    \item product title;
    \item product image URL;
    \item product URL;
    \item star rating;
    \item number of reviews;
    \item price;
    \item list price;
    \item original category ID;
    \item best-seller indicator;
    \item recently purchased count.
\end{itemize}

The product dataset was merged with a category mapping dataset containing
248 original category IDs. All product records were successfully associated
with a category in the merged dataset.

For the language-model classification task, the original categories were
mapped into the following 16 higher-level main categories:

\begin{enumerate}
    \item Arts, Crafts \& Party Supplies
    \item Automotive
    \item Baby Products
    \item Beauty \& Personal Care
    \item Electronics \& Computers
    \item Fashion, Shoes \& Luggage
    \item Gift Cards
    \item Health \& Household
    \item Home \& Kitchen
    \item Industrial \& Scientific
    \item Pet Supplies
    \item Smart Home
    \item Sports \& Outdoors
    \item Tools \& Home Improvement
    \item Toys \& Games
    \item Video Games
\end{enumerate}


\subsection{Problem Formulation}

Let a product title or user request be represented by text input $x$.

The classification task is to predict one category

\[
\hat{y} = f_{\theta}(x),
\]

where

\[
\hat{y} \in \mathcal{C},
\]

and $\mathcal{C}$ is the predefined set of 16 main categories.

Unlike open-ended text generation, this is a closed-set, single-label,
multiclass classification problem.

A prediction is considered correct only when

\[
\hat{y} = y,
\]

where $y$ is the ground-truth category from the dataset taxonomy.

This distinction is important because a general-purpose model may generate a
semantically reasonable but non-matching category. For example,

\begin{quote}
Ground truth: \texttt{Electronics \& Computers}

Model output: \texttt{Computer Accessories}
\end{quote}

Although the output may be semantically related, it does not belong to the
predefined label set and therefore cannot directly be used for deterministic
database filtering.


\subsection{Prompt Engineering versus Fine-Tuning}

Before fine-tuning, prompt engineering was used to improve the behavior of
the base models.

Instead of asking the model to freely classify a product, the prompt
explicitly included all 16 valid labels and instructed the model to return
exactly one category from the provided list.

For example:

\begin{quote}
``Classify the following product into exactly one of these 16 categories.
Only return the category name.''
\end{quote}

This significantly reduces unconstrained category generation. However,
prompt engineering alone does not modify the model parameters.

Therefore, it can specify the desired output format but does not explicitly
teach the model the dataset-specific boundaries between categories.

For example, the model may know that a gaming mouse is related to
``computer accessories,'' while the project taxonomy specifically requires
the label ``Electronics \& Computers.''

Fine-tuning addresses this issue by repeatedly exposing the model to
training examples of the form

\[
\text{Product Title} \rightarrow
\text{Predefined Main Category}.
\]

The model therefore learns both the output format and the specific category
mapping used by our dataset.


% =========================================================
\section{Design}

\subsection{Overall System Architecture}

The proposed shopping assistant follows the pipeline:

\begin{center}
\fbox{
\begin{minipage}{0.92\linewidth}
\centering
User Query
$\rightarrow$
Structured Query Extraction
$\rightarrow$
Fine-Tuned Category Classifier
$\rightarrow$
SQL Filtering
$\rightarrow$
Semantic Vector Retrieval
$\rightarrow$
Recommendation Generation
\end{minipage}}
\end{center}

A user may provide a query such as:

\begin{quote}
``I want a wireless gaming mouse under \$50 with at least four stars.''
\end{quote}

The structured query extraction component identifies constraints such as:

\begin{itemize}
    \item product description: wireless gaming mouse;
    \item maximum price: \$50;
    \item minimum star rating: 4;
    \item brand: optional;
    \item minimum review count: optional;
    \item sorting criterion.
\end{itemize}

The fine-tuned model independently predicts:

\begin{quote}
\texttt{Electronics \& Computers}
\end{quote}

This predicted category is then added as a strict database filter.

The database query therefore combines symbolic constraints with semantic
information rather than relying only on natural-language generation.


\subsection{Structured Database Retrieval}

Structured properties such as price, category, star rating, review count,
and brand can be represented as SQL filters.

A simplified query may take the form:

\begin{verbatim}
SELECT *
FROM amazon_products
WHERE main_category = 'Electronics & Computers'
  AND price <= 50
  AND stars >= 4;
\end{verbatim}

This prevents the recommendation model from inventing product attributes
that do not exist in the database.


\subsection{Semantic Retrieval}

Exact keyword matching alone may be too restrictive. For example,
``portable gaming mouse'' and ``wireless lightweight mouse'' may represent
similar user intent without containing exactly the same keywords.

To improve semantic matching, product-title embeddings are stored in a
vector representation.

The project uses the embedding model

\begin{quote}
\texttt{BAAI/bge-base-en-v1.5}
\end{quote}

with an embedding dimension of 768.

Semantic similarity can then be used to rank candidate products after
structured filters have removed invalid candidates.


% =========================================================
\section{Methodology}

\subsection{Model Selection}

Three lightweight model families were selected:

\begin{itemize}
    \item Qwen3.5-0.8B;
    \item Llama 3.2-1B;
    \item Gemma 3 1B IT.
\end{itemize}

The models were selected for three main reasons.

First, all three are approximately within the same small-model parameter
range, making comparison more meaningful than comparing a 1B model directly
against a much larger model.

Second, they represent different model families. Comparing different
families provides more information than comparing several checkpoints from
the same architecture.

Third, their size allows fine-tuning and inference using the computational
resources available for this project.


\subsection{Why Approximately 1B Parameters?}

The category prediction task is significantly narrower than general
open-domain language generation.

The model does not need to:

\begin{itemize}
    \item generate long essays;
    \item perform complex multi-step reasoning;
    \item memorize the entire Amazon catalog;
    \item answer arbitrary world-knowledge questions.
\end{itemize}

Instead, it must map a short product title or shopping request to one of only
16 possible output labels.

For such a constrained problem, using a much larger model would substantially
increase training memory, inference latency, and computational cost without
necessarily providing a proportional improvement in classification
performance.

Small models also make local deployment possible and reduce dependence on
external API calls.

Our hypothesis is therefore that task-specific fine-tuning of a small model
can provide better cost-performance trade-offs than using an unfine-tuned
large model for a narrow classification task.

Larger models remain a possible extension and are discussed in
Section~\ref{sec:discussion}.


\subsection{Supervised Fine-Tuning}

Each training example contains a product title and its ground-truth
main category.

Conceptually, each example follows:

\begin{quote}
Instruction: Classify the product into one main category.

Input: Product title

Output: One of the 16 predefined category labels.
\end{quote}

The model is trained using next-token prediction over the supervised
instruction-response sequence.


\subsection{LoRA Fine-Tuning of Gemma 3 1B IT}

Gemma 3 1B IT was fine-tuned using Low-Rank Adaptation (LoRA).

Instead of updating every parameter in the approximately 1B-parameter model,
LoRA freezes the original model weights and introduces trainable low-rank
matrices into selected linear projection layers.

For an original weight matrix $W$, LoRA represents the adapted weight as

\[
W' = W + BA,
\]

where $A$ and $B$ are low-rank trainable matrices.

The rank $r$ is much smaller than the original matrix dimensions, allowing
the model to be adapted using significantly fewer trainable parameters.

For the Gemma experiment, the LoRA configuration was:

\begin{table}[h]
\caption{Gemma 3 1B IT LoRA configuration.}
\label{tab:gemma-config}
\centering
\begin{tabular}{ll}
\toprule
Parameter & Value \\
\midrule
LoRA rank $r$ & 16 \\
LoRA alpha & 32 \\
LoRA dropout & 0.05 \\
Learning rate & $2\times10^{-4}$ \\
Training epochs & 1 \\
Per-device batch size & 4 \\
Gradient accumulation steps & 4 \\
Effective batch size & 16 \\
Maximum sequence length & 512 \\
Training precision & BF16 \\
Gradient checkpointing & Disabled \\
\bottomrule
\end{tabular}
\end{table}

LoRA adapters were applied to:

\begin{itemize}
    \item \texttt{q\_proj};
    \item \texttt{k\_proj};
    \item \texttt{v\_proj};
    \item \texttt{o\_proj};
    \item \texttt{gate\_proj};
    \item \texttt{up\_proj};
    \item \texttt{down\_proj}.
\end{itemize}

These layers cover both the attention projections and MLP projection layers.

The implementation used the Hugging Face Transformers, PEFT, and TRL
libraries.

The supervised fine-tuning procedure can be summarized as:

\begin{enumerate}
    \item Load the pretrained instruction-tuned language model.
    \item Freeze the original model parameters.
    \item Insert LoRA adapters into selected projection layers.
    \item Format product titles and labels as instruction-response pairs.
    \item Optimize only the LoRA parameters using supervised fine-tuning.
    \item Save the LoRA adapter weights.
    \item Load the adapter during inference and generate exactly one category.
\end{enumerate}


\subsection{Training Data}

The complete balanced classification dataset contained approximately
71,586 examples before the final experimental sampling procedure.

For the reported Gemma training run:

\begin{itemize}
    \item training examples: 10,000;
    \item testing examples: 1,000.
\end{itemize}

\todo{Confirm whether Qwen and Llama used exactly the same 10,000/1,000
samples. If not, report the exact sample counts for each model here.}

An exact-title leakage check found only four overlapping titles between the
sampled training and test sets, corresponding to approximately 0.4\% of the
test titles. This suggests that the measured performance was not primarily
caused by exact title memorization.


% =========================================================
\section{Numerical Experiments}

\subsection{Experimental Setup}

The main experiment compares the three lightweight language models on the
same 16-class product classification task.

The primary objective is not merely to measure whether the generated text is
semantically related to the product. The output must match the predefined
taxonomy so that it can be directly used by the downstream database query.

The following evaluation metrics are considered:

\begin{itemize}
    \item Accuracy;
    \item Macro Precision;
    \item Macro Recall;
    \item Macro F1;
    \item Weighted F1;
    \item Balanced Accuracy;
    \item Exact Label Compliance Rate;
    \item Per-class Precision, Recall, and F1;
    \item Confusion Matrix.
\end{itemize}


\subsection{Evaluation Metrics}

For $N$ samples, classification accuracy is

\[
\mathrm{Accuracy}
=
\frac{1}{N}
\sum_{i=1}^{N}
\mathbb{1}(\hat{y}_i=y_i).
\]

Accuracy provides an intuitive measure of overall correctness.

However, accuracy alone can hide poor performance on categories with fewer
examples. Therefore, Precision, Recall, and F1 are also evaluated separately
for each class.

For class $c$,

\[
\mathrm{Precision}_c =
\frac{TP_c}{TP_c + FP_c},
\]

\[
\mathrm{Recall}_c =
\frac{TP_c}{TP_c + FN_c},
\]

and

\[
F1_c =
2
\frac{
\mathrm{Precision}_c \cdot \mathrm{Recall}_c
}{
\mathrm{Precision}_c + \mathrm{Recall}_c
}.
\]

Macro F1 assigns equal importance to every category:

\[
\mathrm{MacroF1}
=
\frac{1}{C}
\sum_{c=1}^{C}F1_c,
\]

where $C=16$.

This metric is particularly useful because a model cannot obtain a strong
Macro F1 score simply by performing well on the most frequent categories.

Weighted F1 is

\[
\mathrm{WeightedF1}
=
\sum_{c=1}^{C}
\frac{N_c}{N}F1_c.
\]

It measures overall F1 performance while accounting for the class
distribution.

Balanced Accuracy is equivalent to averaging class-wise recall for a
multiclass classification task and is included to evaluate whether smaller
classes are being ignored.

We additionally define \textbf{Exact Label Compliance Rate}:

\[
\mathrm{Compliance}
=
\frac{
\#\{\hat{y}_i \in \mathcal{C}\}
}{
N
}.
\]

This metric is important for an LLM-based classifier because a model can
generate a semantically plausible category that is not one of the 16 allowed
labels.


\subsection{Prompt Engineering Baseline}

Before fine-tuning, we evaluated a prompt-engineering approach in which the
model was explicitly shown all 16 valid category names.

The prompt instructed the model to:

\begin{enumerate}
    \item select exactly one category;
    \item use only a category from the predefined label list;
    \item return only the category name without additional explanation.
\end{enumerate}

This baseline is more appropriate than unconstrained zero-shot generation
because a completely unconstrained model may invent its own taxonomy.

For Gemma 3 1B IT, the prompt-based zero-shot baseline achieved:

\begin{table}[h]
\caption{Prompt-based zero-shot baseline for Gemma 3 1B IT.}
\centering
\begin{tabular}{lc}
\toprule
Metric & Result \\
\midrule
Accuracy & 34.0\% \\
Macro F1 & 13.36\% \\
Weighted F1 & 31.30\% \\
Macro Precision & \todo{calculate} \\
Macro Recall & \todo{calculate} \\
Balanced Accuracy & \todo{calculate} \\
Exact Label Compliance & \todo{calculate} \\
\bottomrule
\end{tabular}
\end{table}

The relatively low performance demonstrates that prompt constraints alone
are insufficient for learning the project-specific taxonomy.

Prompt engineering can tell the model which outputs are allowed, but it does
not teach the model the decision boundary between those categories.


\subsection{Fine-Tuned Model Comparison}

Table~\ref{tab:model-results} compares the three fine-tuned lightweight
models.

\begin{table}[h]
\caption{Fine-tuned model comparison.}
\label{tab:model-results}
\centering
\resizebox{\textwidth}{!}{
\begin{tabular}{lccccccc}
\toprule
Model &
Accuracy &
Macro Precision &
Macro Recall &
Macro F1 &
Weighted F1 &
Balanced Acc. &
Label Compliance \\
\midrule

Qwen3.5-0.8B
& 75.3\%
& \todo{value}
& \todo{value}
& 75.7\%
& 75.3\%
& \todo{value}
& \todo{value}
\\

Llama 3.2-1B
& 87.7\%
& \todo{value}
& \todo{value}
& 81.2\%
& \todo{value}
& \todo{value}
& \todo{value}
\\

Gemma 3 1B IT
& \textbf{93.5\%}
& \todo{value}
& \todo{value}
& \textbf{87.4\%}
& \textbf{93.5\%}
& \todo{value}
& \todo{value}
\\

\bottomrule
\end{tabular}}
\end{table}

Gemma achieved the best observed overall performance, with 93.5\% accuracy
and a Macro F1 of approximately 87.4\%.

The difference between Accuracy and Macro F1 indicates that some classes
remain more difficult than others. Therefore, per-class results were also
examined rather than relying only on aggregate accuracy.


\subsection{Gemma Per-Class Performance}

The fine-tuned Gemma model achieved strong F1 scores for most categories.

Examples include:

\begin{itemize}
    \item Gift Cards: F1 = 1.00;
    \item Industrial \& Scientific: F1 = 0.98;
    \item Pet Supplies: F1 = 0.98;
    \item Video Games: F1 = 0.98;
    \item Electronics \& Computers: F1 = 0.96;
    \item Automotive: F1 = 0.96.
\end{itemize}

The more difficult categories included:

\begin{itemize}
    \item Smart Home: F1 = 0.75;
    \item Sports \& Outdoors: F1 = 0.82;
    \item Health \& Household: F1 = 0.88.
\end{itemize}

In particular, Smart Home achieved high precision but substantially lower
recall, suggesting that the model is conservative when assigning products to
this category.

\todo{Insert confusion matrix figure generated from the final Gemma test set.}

\begin{figure}[h]
  \centering
  \fbox{
  \rule[-.5cm]{0cm}{5.0cm}
  \rule[-.5cm]{8.0cm}{0cm}
  }
  \caption{Confusion matrix of the fine-tuned Gemma 3 1B IT classifier.
  Replace this placeholder with the final confusion-matrix image.}
  \label{fig:confusion}
\end{figure}


\subsection{Training Curves}

\todo{Export final training-loss and validation-loss curves for each model.}

\begin{figure}[h]
  \centering
  \fbox{
  \rule[-.5cm]{0cm}{4.5cm}
  \rule[-.5cm]{8cm}{0cm}
  }
  \caption{Training loss of the fine-tuned models.}
  \label{fig:training-loss}
\end{figure}


\subsection{Efficiency Comparison}

Classification quality is only one factor in model selection.

Because the system is intended to use a lightweight classifier, model
efficiency should also be considered.

\begin{table}[h]
\caption{Model efficiency comparison.}
\label{tab:efficiency}
\centering
\begin{tabular}{lcccc}
\toprule
Model &
Parameters &
Trainable LoRA Params. &
Peak VRAM &
Avg. Inference Time \\
\midrule
Qwen3.5-0.8B
& $\sim$0.8B
& \todo{value}
& \todo{GB}
& \todo{ms/sample} \\

Llama 3.2-1B
& $\sim$1B
& \todo{value}
& \todo{GB}
& \todo{ms/sample} \\

Gemma 3 1B IT
& $\sim$1B
& \todo{value}
& \todo{GB}
& \todo{ms/sample} \\
\bottomrule
\end{tabular}
\end{table}

Additional useful engineering metrics include:

\begin{itemize}
    \item training time;
    \item GPU memory usage;
    \item inference latency;
    \item throughput in samples per second;
    \item number and percentage of trainable parameters.
\end{itemize}

These measurements allow comparison not only in terms of predictive
performance but also deployment efficiency.


\subsection{End-to-End Demo Test}

The final pipeline was also tested using natural-language shopping requests.

For example:

\begin{quote}
``I want a wireless gaming mouse under \$50.''
\end{quote}

The structured query component extracts the price and product-description
constraints, while the fine-tuned Gemma classifier predicts:

\begin{quote}
\texttt{Electronics \& Computers}
\end{quote}

The database then retrieves only candidate products satisfying the category
and structured conditions.

\todo{Insert 2--5 representative demo queries and returned products.}

For each demo query, the report should record:

\begin{itemize}
    \item original user request;
    \item structured query;
    \item predicted main category;
    \item top retrieved products;
    \item whether all hard constraints were satisfied.
\end{itemize}


% =========================================================
\section{Discussion}
\label{sec:discussion}

\subsection{Why Not Use ChatGPT Directly?}

A natural question is why the complete shopping assistant is not implemented
by directly asking a powerful general-purpose model such as ChatGPT to
recommend products.

The main reason is that the objectives of a generative chatbot and a
database retrieval system are different.

A generative model is optimized to generate text that is relevant and useful
in a broad semantic sense. It is not inherently constrained to return only
database records satisfying strict application-specific conditions.

In our exploratory tests, asking for a product such as a mouse below a
specified price could occasionally produce related products or accessories,
such as mouse pads. Although these results may be semantically related to the
request, they are incorrect for a strict product retrieval task.

This becomes more important when multiple hard constraints are present:

\begin{itemize}
    \item product category;
    \item maximum price;
    \item minimum rating;
    \item minimum review count;
    \item brand;
    \item availability in the project database.
\end{itemize}

A database system can guarantee such constraints directly, whereas a
generative LLM cannot guarantee them through prompting alone.

There are also practical considerations.

Using an external large-model API for every classification and retrieval
decision introduces:

\begin{itemize}
    \item API cost;
    \item network latency;
    \item dependence on an external service;
    \item reduced reproducibility;
    \item less control over future model changes;
    \item difficulty guaranteeing exact output schemas.
\end{itemize}

Therefore, our system uses LLMs only where language understanding is useful
and uses deterministic database operations where strict constraints are
required.


\subsection{Why Fine-Tuning is Necessary}

Prompt engineering improved the behavior of the base models by explicitly
providing the 16 allowed labels.

However, a prompt mainly specifies \emph{what the model should do}.
It does not directly teach the model \emph{how our dataset defines each
category}.

For instance, a general language model may reasonably classify a wireless
mouse as:

\begin{quote}
\texttt{Computer Accessories}
\end{quote}

while the project taxonomy requires:

\begin{quote}
\texttt{Electronics \& Computers}.
\end{quote}

The first answer is linguistically reasonable but unusable as a direct
database category filter.

Fine-tuning repeatedly exposes the model to examples from the actual
taxonomy and therefore modifies the model behavior toward the application's
specific decision boundaries.

This difference can be summarized as:

\begin{quote}
\textbf{Prompt engineering:}
``Choose one of these 16 labels.''

\medskip

\textbf{Fine-tuning:}
``Learn what each of these 16 labels means in our dataset.''
\end{quote}

This is supported by the Gemma experiment, where the constrained
prompt-based baseline achieved approximately 34.0\% accuracy, while the
fine-tuned model achieved 93.5\% accuracy.


\subsection{Why Use Small Models Instead of Larger LLMs?}

Larger models generally provide greater general-purpose language and
reasoning capability. However, model size should be matched to the
complexity of the target task.

Our category classifier has:

\begin{itemize}
    \item short textual inputs;
    \item only 16 output classes;
    \item no requirement for long-form generation;
    \item no requirement for broad world knowledge;
    \item no need to memorize product records.
\end{itemize}

Therefore, much of the capability provided by a much larger model is not
required.

A smaller approximately 1B-parameter model provides several advantages:

\begin{enumerate}
    \item \textbf{Lower training cost.}
    LoRA fine-tuning can be performed using limited GPU resources.

    \item \textbf{Lower inference latency.}
    A smaller classifier can make category predictions faster.

    \item \textbf{Lower memory requirement.}
    Small models are easier to deploy locally.

    \item \textbf{Reduced API dependence.}
    The classifier can run without sending every request to an external
    large-model API.

    \item \textbf{Task specialization.}
    Fine-tuning allows a small model to focus its capacity on the narrow
    classification problem.

    \item \textbf{Fair experimental comparison.}
    Models from different families can be compared within a similar
    parameter range.
\end{enumerate}

The 93.5\% Gemma accuracy suggests that a lightweight model has sufficient
capacity for this particular classification task.

However, this does not establish that 1B is universally optimal.
A larger model may improve difficult or ambiguous categories.
A future experiment could compare 1B, 3B, 4B, and larger models while
measuring the trade-off between classification quality, latency, memory,
and training cost.


\subsection{Limitations}

Several limitations remain.

First, the classifier depends primarily on product-title information.
Some products may be difficult to classify without richer descriptions.

Second, some high-level categories have semantic overlap. For example,
Smart Home products may also appear similar to Electronics \& Computers or
Home \& Kitchen.

Third, the dataset does not contain sufficient user interaction history for
a complete collaborative recommendation system. Therefore, the current
system focuses primarily on query-based retrieval rather than long-term
personalized recommendation.

Fourth, evaluation of the complete shopping assistant should include
retrieval quality in addition to category-classification accuracy.


\subsection{Future Retrieval Metrics}

For the final retrieval component, future evaluation should use ranking
metrics such as:

\begin{itemize}
    \item Recall@$K$;
    \item Precision@$K$;
    \item Mean Reciprocal Rank (MRR);
    \item Normalized Discounted Cumulative Gain (NDCG@$K$);
    \item hard-constraint satisfaction rate.
\end{itemize}

These metrics measure whether relevant products appear near the top of the
retrieved list, which cannot be evaluated using classification accuracy
alone.


% =========================================================
\section{Conclusions}

This project investigated whether lightweight language models can be
specialized for product category classification and integrated into a
database-driven shopping assistant.

Three model families---Qwen3.5-0.8B, Llama 3.2-1B, and Gemma 3 1B IT---were
fine-tuned and evaluated on a 16-class Amazon product taxonomy.

Among the tested models, Gemma 3 1B IT achieved the best observed
performance, with 93.5\% classification accuracy, approximately 87.4\%
Macro F1, and approximately 93.5\% Weighted F1.

The experiments also demonstrate the importance of task-specific
fine-tuning. Prompt engineering can restrict the expected output format, but
it does not fully teach a general-purpose model the category boundaries used
by a specific dataset. Fine-tuning substantially improved the model's
ability to map product titles to the required closed-set taxonomy.

The broader system design also suggests that an effective shopping assistant
does not require one large language model to perform every task. Instead,
different components can be used according to their strengths: language
models for natural-language understanding, fine-tuned small models for
controlled classification, SQL for hard constraints, vector embeddings for
semantic similarity, and a generative model for final user-facing
explanations.

Future work will evaluate the semantic retrieval component using ranking
metrics such as Recall@$K$, MRR, and NDCG, investigate richer product
descriptions, and compare the performance-cost trade-off of larger language
models.


% =========================================================
\section*{References}

\small

[1] Hu, E. J., Shen, Y., Wallis, P., Allen-Zhu, Z., Li, Y.,
Wang, S., Wang, L., and Chen, W. (2021).
LoRA: Low-Rank Adaptation of Large Language Models.
\textit{arXiv preprint arXiv:2106.09685}.

\medskip

[2] \todo{Add official Gemma 3 technical report/model card reference.}

\medskip

[3] \todo{Add official Llama 3.2 model reference.}

\medskip

[4] \todo{Add official Qwen model/model-card reference matching the exact
checkpoint used in the project.}

\medskip

[5] \todo{Add BAAI BGE embedding model reference.}

\medskip

[6] \todo{Add Amazon dataset/Kaggle source reference.}

\medskip

[7] \todo{Add Hugging Face Transformers reference.}

\medskip

[8] \todo{Add PEFT/TRL references if required.}


% =========================================================
\section*{Appendix}

\subsection*{A. Final Hyperparameter Configuration}

\todo{Add complete Qwen and Llama hyperparameter tables so that all
three experiments can be reproduced.}

\subsection*{B. Full Classification Reports}

\todo{Insert the full per-category Precision, Recall, F1, and support
tables for Qwen, Llama, and Gemma.}

\subsection*{C. Example Queries}

\todo{Add additional end-to-end shopping assistant examples.}

\end{document}